% Options for packages loaded elsewhere
% Options for packages loaded elsewhere
\PassOptionsToPackage{unicode}{hyperref}
\PassOptionsToPackage{hyphens}{url}
\PassOptionsToPackage{dvipsnames,svgnames,x11names}{xcolor}
%
\documentclass[journal=iecr,manuscript=]{achemso}
\usepackage[version=3]{mhchem}
\newcommand*\mycommand[1]{\texttt{\emph{#1}}}


\usepackage{xcolor}
\usepackage{amsmath,amssymb}
\setcounter{secnumdepth}{-\maxdimen} % remove section numbering
\usepackage{iftex}
\ifPDFTeX
  \usepackage[T1]{fontenc}
  \usepackage[utf8]{inputenc}
  \usepackage{textcomp} % provide euro and other symbols
\else % if luatex or xetex
  \usepackage{unicode-math} % this also loads fontspec
  \defaultfontfeatures{Scale=MatchLowercase}
  \defaultfontfeatures[\rmfamily]{Ligatures=TeX,Scale=1}
\fi
\usepackage{lmodern}
\ifPDFTeX\else
  % xetex/luatex font selection
\fi
% Use upquote if available, for straight quotes in verbatim environments
\IfFileExists{upquote.sty}{\usepackage{upquote}}{}
\IfFileExists{microtype.sty}{% use microtype if available
  \usepackage[]{microtype}
  \UseMicrotypeSet[protrusion]{basicmath} % disable protrusion for tt fonts
}{}
\makeatletter
\@ifundefined{KOMAClassName}{% if non-KOMA class
  \IfFileExists{parskip.sty}{%
    \usepackage{parskip}
  }{% else
    \setlength{\parindent}{0pt}
    \setlength{\parskip}{6pt plus 2pt minus 1pt}}
}{% if KOMA class
  \KOMAoptions{parskip=half}}
\makeatother
% Make \paragraph and \subparagraph free-standing
\makeatletter
\ifx\paragraph\undefined\else
  \let\oldparagraph\paragraph
  \renewcommand{\paragraph}{
    \@ifstar
      \xxxParagraphStar
      \xxxParagraphNoStar
  }
  \newcommand{\xxxParagraphStar}[1]{\oldparagraph*{#1}\mbox{}}
  \newcommand{\xxxParagraphNoStar}[1]{\oldparagraph{#1}\mbox{}}
\fi
\ifx\subparagraph\undefined\else
  \let\oldsubparagraph\subparagraph
  \renewcommand{\subparagraph}{
    \@ifstar
      \xxxSubParagraphStar
      \xxxSubParagraphNoStar
  }
  \newcommand{\xxxSubParagraphStar}[1]{\oldsubparagraph*{#1}\mbox{}}
  \newcommand{\xxxSubParagraphNoStar}[1]{\oldsubparagraph{#1}\mbox{}}
\fi
\makeatother

\usepackage{color}
\usepackage{fancyvrb}
\newcommand{\VerbBar}{|}
\newcommand{\VERB}{\Verb[commandchars=\\\{\}]}
\DefineVerbatimEnvironment{Highlighting}{Verbatim}{commandchars=\\\{\}}
% Add ',fontsize=\small' for more characters per line
\usepackage{framed}
\definecolor{shadecolor}{RGB}{241,243,245}
\newenvironment{Shaded}{\begin{snugshade}}{\end{snugshade}}
\newcommand{\AlertTok}[1]{\textcolor[rgb]{0.68,0.00,0.00}{#1}}
\newcommand{\AnnotationTok}[1]{\textcolor[rgb]{0.37,0.37,0.37}{#1}}
\newcommand{\AttributeTok}[1]{\textcolor[rgb]{0.40,0.45,0.13}{#1}}
\newcommand{\BaseNTok}[1]{\textcolor[rgb]{0.68,0.00,0.00}{#1}}
\newcommand{\BuiltInTok}[1]{\textcolor[rgb]{0.00,0.23,0.31}{#1}}
\newcommand{\CharTok}[1]{\textcolor[rgb]{0.13,0.47,0.30}{#1}}
\newcommand{\CommentTok}[1]{\textcolor[rgb]{0.37,0.37,0.37}{#1}}
\newcommand{\CommentVarTok}[1]{\textcolor[rgb]{0.37,0.37,0.37}{\textit{#1}}}
\newcommand{\ConstantTok}[1]{\textcolor[rgb]{0.56,0.35,0.01}{#1}}
\newcommand{\ControlFlowTok}[1]{\textcolor[rgb]{0.00,0.23,0.31}{\textbf{#1}}}
\newcommand{\DataTypeTok}[1]{\textcolor[rgb]{0.68,0.00,0.00}{#1}}
\newcommand{\DecValTok}[1]{\textcolor[rgb]{0.68,0.00,0.00}{#1}}
\newcommand{\DocumentationTok}[1]{\textcolor[rgb]{0.37,0.37,0.37}{\textit{#1}}}
\newcommand{\ErrorTok}[1]{\textcolor[rgb]{0.68,0.00,0.00}{#1}}
\newcommand{\ExtensionTok}[1]{\textcolor[rgb]{0.00,0.23,0.31}{#1}}
\newcommand{\FloatTok}[1]{\textcolor[rgb]{0.68,0.00,0.00}{#1}}
\newcommand{\FunctionTok}[1]{\textcolor[rgb]{0.28,0.35,0.67}{#1}}
\newcommand{\ImportTok}[1]{\textcolor[rgb]{0.00,0.46,0.62}{#1}}
\newcommand{\InformationTok}[1]{\textcolor[rgb]{0.37,0.37,0.37}{#1}}
\newcommand{\KeywordTok}[1]{\textcolor[rgb]{0.00,0.23,0.31}{\textbf{#1}}}
\newcommand{\NormalTok}[1]{\textcolor[rgb]{0.00,0.23,0.31}{#1}}
\newcommand{\OperatorTok}[1]{\textcolor[rgb]{0.37,0.37,0.37}{#1}}
\newcommand{\OtherTok}[1]{\textcolor[rgb]{0.00,0.23,0.31}{#1}}
\newcommand{\PreprocessorTok}[1]{\textcolor[rgb]{0.68,0.00,0.00}{#1}}
\newcommand{\RegionMarkerTok}[1]{\textcolor[rgb]{0.00,0.23,0.31}{#1}}
\newcommand{\SpecialCharTok}[1]{\textcolor[rgb]{0.37,0.37,0.37}{#1}}
\newcommand{\SpecialStringTok}[1]{\textcolor[rgb]{0.13,0.47,0.30}{#1}}
\newcommand{\StringTok}[1]{\textcolor[rgb]{0.13,0.47,0.30}{#1}}
\newcommand{\VariableTok}[1]{\textcolor[rgb]{0.07,0.07,0.07}{#1}}
\newcommand{\VerbatimStringTok}[1]{\textcolor[rgb]{0.13,0.47,0.30}{#1}}
\newcommand{\WarningTok}[1]{\textcolor[rgb]{0.37,0.37,0.37}{\textit{#1}}}

\usepackage{longtable,booktabs,array}
\usepackage{calc} % for calculating minipage widths
% Correct order of tables after \paragraph or \subparagraph
\usepackage{etoolbox}
\makeatletter
\patchcmd\longtable{\par}{\if@noskipsec\mbox{}\fi\par}{}{}
\makeatother
% Allow footnotes in longtable head/foot
\IfFileExists{footnotehyper.sty}{\usepackage{footnotehyper}}{\usepackage{footnote}}
\makesavenoteenv{longtable}
\usepackage{graphicx}
\makeatletter
\newsavebox\pandoc@box
\newcommand*\pandocbounded[1]{% scales image to fit in text height/width
  \sbox\pandoc@box{#1}%
  \Gscale@div\@tempa{\textheight}{\dimexpr\ht\pandoc@box+\dp\pandoc@box\relax}%
  \Gscale@div\@tempb{\linewidth}{\wd\pandoc@box}%
  \ifdim\@tempb\p@<\@tempa\p@\let\@tempa\@tempb\fi% select the smaller of both
  \ifdim\@tempa\p@<\p@\scalebox{\@tempa}{\usebox\pandoc@box}%
  \else\usebox{\pandoc@box}%
  \fi%
}
% Set default figure placement to htbp
\def\fps@figure{htbp}
\makeatother





\setlength{\emergencystretch}{3em} % prevent overfull lines

\providecommand{\tightlist}{%
  \setlength{\itemsep}{0pt}\setlength{\parskip}{0pt}}





\makeatletter
\@ifpackageloaded{caption}{}{\usepackage{caption}}
\AtBeginDocument{%
\ifdefined\contentsname
  \renewcommand*\contentsname{Table of contents}
\else
  \newcommand\contentsname{Table of contents}
\fi
\ifdefined\listfigurename
  \renewcommand*\listfigurename{List of Figures}
\else
  \newcommand\listfigurename{List of Figures}
\fi
\ifdefined\listtablename
  \renewcommand*\listtablename{List of Tables}
\else
  \newcommand\listtablename{List of Tables}
\fi
\ifdefined\figurename
  \renewcommand*\figurename{Figure}
\else
  \newcommand\figurename{Figure}
\fi
\ifdefined\tablename
  \renewcommand*\tablename{Table}
\else
  \newcommand\tablename{Table}
\fi
}
\@ifpackageloaded{float}{}{\usepackage{float}}
\floatstyle{ruled}
\@ifundefined{c@chapter}{\newfloat{codelisting}{h}{lop}}{\newfloat{codelisting}{h}{lop}[chapter]}
\floatname{codelisting}{Listing}
\newcommand*\listoflistings{\listof{codelisting}{List of Listings}}
\makeatother
\makeatletter
\makeatother
\makeatletter
\@ifpackageloaded{caption}{}{\usepackage{caption}}
\@ifpackageloaded{subcaption}{}{\usepackage{subcaption}}
\makeatother
\usepackage{bookmark}
\IfFileExists{xurl.sty}{\usepackage{xurl}}{} % add URL line breaks if available
\urlstyle{same}
\hypersetup{
  pdftitle={Insights into the Optimal Flow Rule for GSHP Systems},
  pdfauthor={Your Name},
  pdfkeywords={Ground Source Heat Pump, Flow Control, Thermal-Hydraulic
Coupling, Optimization},
  colorlinks=true,
  linkcolor={blue},
  filecolor={Maroon},
  citecolor={Blue},
  urlcolor={Blue},
  pdfcreator={LaTeX via pandoc}}


\author{Your Name}
\affiliation{ Your Institution,  }


\email{your.email@institution.edu}


\keywords{Ground Source Heat PumpFlow ControlThermal-Hydraulic
CouplingOptimization}

\title[]{Insights into the Optimal Flow Rule for GSHP Systems}
\makeatletter
\begin{document}
\maketitle
\begin{abstract}
This paper investigates the question of flow control in ground source
heat pump (GSHP) systems. We focus on systems with a single heat pump
connected hydraulically to a borehole without decoupling tanks. The
study develops a coupled thermal-hydraulic model that considers
steady-state conditions with the borehole wall as the boundary
condition. The model incorporates principles of fluid mechanics and heat
transfer while accounting for thermal properties of both ground and heat
pump components.
\end{abstract}


\section{Abstract:}\label{abstract}

This paper will look into the question of the flow control in GSHP
systems. We will focus on systems with a single heat pump connected
hydraulically to a borehole in terms of hydraulics, no decoupling tank.

\section{Introduction}\label{introduction}

Heat pump in genreal and GSHP in particular are essential tools allowing
for the decarbonization of the heating sector. The question of the flow
control in GSHP systems is a central one, as it allows for the
optimization of the system performance and the reduction of energy
consumption. The cost of changing a few variables in a control algorithm
is minimal. Therefore, it is important to understand the hydraulic
aspects of GSHP systems in order to design effective control strategies.
The flow control in GSHP systems is a complex problem that involves the
interaction between the heat pump, the borehole, and the ground. The
flow rate through the borehole is influenced by several factors,
including the temperature difference between the fluid entering and
leaving the borehole, the thermal conductivity of the ground, and the
pressure drop across the borehole. In addition, the flow rate is also
affected by the control strategy used to operate the heat pump.

In this paper, we will focus on the hydraulic aspects of GSHP systems,
specifically on the flow control in systems with a single heat pump
connected hydraulically to a borehole. The model developped considers a
steady state model of the inner problem, and sets the borehole wall as
the boundary condition. The model is implemented in ModelingToolkit.jl,
a Julia package for modeling and simulation of dynamic systems. The
model is based on the principles of fluid mechanics and heat transfer,
and takes into account the thermal properties of the ground and the heat
pump.

\section{Literature:}\label{literature}

The question of the control of HP is not particularly new. Indeed, manyt
authors

\section{Methods:}\label{methods}

The GSHP system model consists of three coupled submodels that are
solved in sequence: thermodynamic cycle (thermal system), borehole heat
exchanger, and pump hydraulics (hydraulic system). The steady-state
system is solved using ModelingToolkit.jl with symbolic manipulation and
NonlinearSolve.jl for robust Newton-Raphson solution methods. The
thermal and hydraulic systems are solved separately but with tight
coupling through shared variables and parameter exchange.

\subsection{System Architecture:}\label{system-architecture}

The model implements a modular approach where the thermal system
determines the thermodynamic cycle performance, heat loads, and required
mass flows, while the hydraulic system calculates pressure drops, pump
power, and validates the hydraulic feasibility of the thermal solution.
The systems are coupled through:

\begin{enumerate}
\def\labelenumi{\arabic{enumi}.}
\tightlist
\item
  \textbf{Mass flow coupling}: The secondary fluid mass flow from the
  thermal system becomes the design flow for the hydraulic system
\item
  \textbf{Temperature coupling}: Borehole inlet/outlet temperatures link
  the thermal model to the hydraulic thermal properties\\
\item
  \textbf{Performance feedback}: Pump power from the hydraulic system
  feeds back into overall system performance calculations
\end{enumerate}

\subsection{Thermodynamic Cycle Model:}\label{thermodynamic-cycle-model}

The thermodynamic cycle implements a comprehensive R134A refrigerant
circuit with the following components:

\subsubsection{Evaporator:}\label{evaporator}

A moving-boundary plate heat exchanger evaporator with dynamically
allocated regions for two-phase evaporation and single-phase
superheating. The model features:

\begin{itemize}
\tightlist
\item
  \textbf{Geometric modeling}: Chevron plate geometry with adjustable
  plate spacing (b), number of plates (np), width (W), enlargement
  factor (ϕ), and chevron angle (β)
\item
  \textbf{Two-zone heat transfer}: The evaporator length is dynamically
  split between evaporation and superheating zones using a smooth
  transition function \texttt{zap(uz\_fr\_sh)} that maps the unbounded
  variable \texttt{uz\_fr\_sh} to the {[}0,1{]} interval
\item
  \textbf{Heat transfer correlations}:

  \begin{itemize}
  \tightlist
  \item
    Two-phase region: Amalfi correlation for convective boiling in plate
    heat exchangers with quality-dependent averaging
  \item
    Superheating region: Martin correlations for single-phase plate heat
    exchanger flow with laminar-turbulent transition
  \end{itemize}
\item
  \textbf{Boundary conditions}: Fixed superheat degree (ΔTSH) and
  secondary fluid temperature difference (ΔTSF) as control variables
\end{itemize}

\subsubsection{Compressor:}\label{compressor}

The scroll compressor model incorporates: - \textbf{Volumetric
displacement}: Fixed displacement volume with frequency-controlled
capacity - \textbf{Isentropic efficiency}: Constant isentropic
efficiency

\subsubsection{Expansion Valve:}\label{expansion-valve}

Modeled as an isenthalpic expansion process. A constraint imposed in the
solver is the superheating degree.

\subsubsection{Condenser:}\label{condenser}

Simplified model assuming constant condensing temperature and specified
subcooling degree (ΔTSC). The detailed condenser heat transfer is not
modeled, focusing on the evaporator side.

\subsection{Borehole Heat Exchanger
Model:}\label{borehole-heat-exchanger-model}

The U-tube grouted borehole employs two sub models:

\subsubsection{Thermal Model:}\label{thermal-model}

\begin{itemize}
\tightlist
\item
  \textbf{Analytical solution}: Uses eigenvalue decomposition of coupled
  1D ODEs to solve the temperature profiles along the borehole depth
\item
  \textbf{Multipole method}: For the grout resistance, we implement
  thermal resistance calculations between pipes and to the borehole wall
  using multipole approximations for complex geometry effects
  \textbf{ref needed}
\item
  \textbf{Heat transfer coefficients}: Derived from the borehole
  geometry and dynamical fluid properties
\item
  \textbf{Boundary conditions}: Uniform borehole wall temperature as the
  thermal boundary condition
\item
  \textbf{Temperature profiles}: Analytical expressions for fluid
  temperatures in both legs as functions of depth
\end{itemize}

\subsubsection{Hydraulic Integration:}\label{hydraulic-integration}

The borehole thermal model is connected to the hydraulic system through
temperature-dependent fluid properties calculated via CoolProp
interpolants registered as symbolic functions in ModelingToolkit.

\subsection{Pump and Hydraulic Model:}\label{pump-and-hydraulic-model}

Note: the hydraulic mode is decoupled from the thermal model, the
solution from the thermal model is used as a boundary condition for the
hydraulic model. There is no feedback from hydraulic losses to the
thermal model, which is a limitation that affects the high flow
simulations (better performances are expected in reality thanks to the
added heat).

\subsubsection{Pump Performance:}\label{pump-performance}

The centrifugal pump follows dimensionless performance curves based on
affinity laws: - \textbf{Head coefficient}:
\(\psi = \psi_0 - K \cdot \varphi^2\) where \(\varphi\) is the flow
coefficient - \textbf{Efficiency}:
\(\eta = \eta_{\text{max}} \left(1 - \left(\frac{\varphi - \varphi_{\text{BEP}}}{\varphi_{\text{BEP}}}\right)^2\right)\)
with parabolic efficiency curve around the best efficiency point (BEP) -
\textbf{Power calculation}: Electrical power including motor efficiency

\subsubsection{Hydraulic Loss
Calculations:}\label{hydraulic-loss-calculations}

\begin{itemize}
\tightlist
\item
  \textbf{Pipe friction}: Darcy-Weisbach equations with smooth
  transition between laminar and turbulent regimes using error function
  blending
\item
  \textbf{Friction factor}: Enhanced Colebrook-White correlations with
  automatic regime detection
\item
  \textbf{Heat exchanger losses}: Chevron-plate correlations based on
  Martin's work for pressure drop in plate heat exchangers
\item
  \textbf{Component-wise pressure drops}: Individual calculation for
  each system component (evaporator, borehole, connecting pipes)
\end{itemize}

\subsection{System Coupling and Solution
Strategy:}\label{system-coupling-and-solution-strategy}

\subsubsection{Sequential Solution
Approach:}\label{sequential-solution-approach}

\begin{enumerate}
\def\labelenumi{\arabic{enumi}.}
\tightlist
\item
  \textbf{Thermal system solution}: Solve the thermodynamic cycle for
  given boundary conditions (borehole wall temperature, condensing
  conditions, control setpoints)
\item
  \textbf{Parameter extraction}: Extract mass flows, temperatures, and
  thermal loads from the thermal solution
\item
  \textbf{Hydraulic system solution}: Solve pump and pressure drop
  calculations using thermal system outputs as boundary conditions
\item
  \textbf{Performance calculation}: Combine thermal and hydraulic
  results to compute overall system performance metrics \#\#\# System
  Coupling and Solution Strategy:
\end{enumerate}

The system architecture follows a sequential solution approach as
illustrated in the system flowchart. The process begins with boundary
conditions (wall temperature, condensing conditions, and control
setpoints) feeding into the thermal system solution. The thermal system
encompasses both the thermodynamic cycle (evaporator, compressor,
condenser, expansion valve) and the borehole heat exchanger thermal
model.

Key coupling points include: - \textbf{Thermal-to-Hydraulic Transfer}:
Mass flow rates and temperatures from the thermal solution become
boundary conditions for the hydraulic system - \textbf{Performance
Integration}: Results from both thermal and hydraulic systems combine to
calculate overall performance metrics (COP1, COP2, COP3) -
\textbf{Auxiliary Heating Assessment}: When heat pump capacity is
insufficient, auxiliary heating requirements are calculated

\emph{Note: A detailed system flowchart will be added in the final
version.}

\subsection{Performance Metrics:}\label{performance-metrics}

The model calculates multiple performance indicators:

\begin{itemize}
\tightlist
\item
  \textbf{Thermal COP}: \(COP_1 = \frac{\dot{Q}_{evap}}{W_{comp}}\)
  (refrigeration cycle only)
\item
  \textbf{System COP}:
  \(COP_2 = \frac{\dot{Q}_{evap}}{W_{comp} + P_{pump}}\) (including pump
  power)
\item
  \textbf{Advanced metrics}: Includes provisions for auxiliary heating
  assessment when heat pump capacity is insufficient
\end{itemize}

\subsection{Numerical Implementation:}\label{numerical-implementation}

\begin{itemize}
\tightlist
\item
  \textbf{Symbolic computation}: Full symbolic derivation of equations
  using ModelingToolkit.jl
\item
  \textbf{Robust solvers}: Multiple nonlinear solver options including
  RobustMultiNewton and TrustRegion methods
\item
  \textbf{Automatic differentiation}: Forward-mode AD for enhanced
  convergence properties
\item
  \textbf{Continuation methods}: Built-in parameter sweep capabilities
  with automatic stepping and convergence control
\end{itemize}

The model implementation prioritizes modularity and extensibility,
allowing for easy modification of individual components while
maintaining system-level coupling consistency. The thermal-hydraulic
separation enables efficient parameter studies and optimization while
ensuring physical consistency through the coupling interface.

\emph{Detailed system parameters and correlations are provided in
Appendix A.}

\section{Results:}\label{results}

\subsection{COP3 vs.~ΔT}\label{cop3-vs.-ux3b4t}

This figure shows the main metric we are considering: COP3 against the
control strategy under investigation: constant temperature difference
across the evaporator/BHE.

There seems to be an optimal

TDL: look into the position of the boundary trace the temperatures in
the evap trace the temperatures in the bhe: get the exponents + ground
temperature reynolds + transitions (illustrer sur les graphs par des
lignes) show limit power pump too UA?

\subsection{Insights into the physical
process}\label{insights-into-the-physical-process}

Position of the boundary within the evaporator The faster the compressor
runs, the more evaporation fraction is used

\subsection{Robustness}\label{robustness}

\subsubsection{Partial load}\label{partial-load}

Chosing the COP3 as a metric allows for the comparison of ``faster
compressor vs.~slower compressor + aux electrical power'' Assuming the
HP can tackle the entirety of the load (P\_aux\_heating = 0.0)

Looking at the graph,

\subsubsection{Change of operating
point}\label{change-of-operating-point}

To establish baseline performance characteristics, we initially examine
a fundamental question: what constitutes the optimal rule-based control
strategy under standard operating conditions?

The test case represents a typical Swedish GSHP installation. The system
features a 250m deep borehole with a 40mm outer diameter polyethylene
single U-tube collector. Complete technical specifications and thermal
properties are detailed in Appendix A.

\subsection{Heatmap of COP2 vs T wall vs T
condenser}\label{heatmap-of-cop2-vs-t-wall-vs-t-condenser}

\subsection{A fine tuned strategy}\label{a-fine-tuned-strategy}

We propose different curves for different operational modes. The
question of the linear transition is central.

\section{Conclusions}\label{conclusions}

It seems that the widespread rule based controll perform decently in the
test case, however, it remains to be seen if the same heuristic does as
well when it comes to eg. coaxial BHE. We suggest that the models be
leveraged in order to produce better heuristics.

\subsection{Nomenclature}\label{nomenclature}

\subsubsection{Symbols}\label{symbols}

\textbf{Geometric Parameters:} - \(D\) = Impeller diameter {[}m{]} -
\(D_{bh}\) = Borehole diameter {[}m{]}\\
- \(D_{pipe,i}\) = Pipe inner diameter {[}m{]} - \(D_{pipe,o}\) = Pipe
outer diameter {[}m{]} - \(H\) = Pump head {[}m{]} - \(L\) = Plate
length {[}m{]} - \(L_b\) = Borehole depth {[}m{]} - \(N_{plates}\) =
Number of plates {[}-{]} - \(W\) = Plate width {[}m{]} - \(b\) = Channel
spacing {[}m{]} - \(s\) = Shank spacing {[}m{]} - \(t_{pipe}\) = Pipe
wall thickness {[}m{]} - \(t_{plate}\) = Plate thickness {[}m{]}

\textbf{Thermophysical Properties:} - \(\lambda\) = Thermal conductivity
{[}W/(m·K){]} - \(\rho\) = Density {[}kg/m³{]} - \(\chi\) = Mass
fraction {[}-{]} - \(\beta\) = Chevron angle {[}rad{]} - \(h\) =
Specific enthalpy {[}kJ/kg{]}

\textbf{Flow and Transport Variables:} - \(\dot{V}\) = Volumetric flow
rate {[}m³/s{]} - \(\dot{m}\) = Mass flow rate {[}kg/s{]} - \(\dot{Q}\)
= Heat transfer rate {[}W{]} - \(N\) = Rotational speed {[}rpm{]} -
\(P\) = Pressure {[}bar, Pa{]} - \(T\) = Temperature {[}K, °C{]} -
\(\Delta T\) = Temperature difference {[}K{]}

\textbf{Dimensionless Parameters:} - \(\phi\) = Flow coefficient,
\(\dot{V}/(N \cdot D^3)\) {[}-{]} - \(\psi\) = Head coefficient,
\(g \cdot H/(N^2 \cdot D^2)\) {[}-{]} - \(\psi_0\) = Shutoff head
coefficient {[}-{]} - \(K\) = Head curve slope parameter {[}-{]}

\textbf{Performance Parameters:} - \(COP\) = Coefficient of performance
{[}-{]} - \(W\) = Work/Power {[}W{]} - \(\eta\) = Efficiency {[}-{]} -
\(f\) = Frequency {[}Hz{]} - \(g\) = Gravitational acceleration
{[}m/s²{]}

\subsubsection{Subscripts}\label{subscripts}

\begin{itemize}
\tightlist
\item
  \(BEP\) = Best efficiency point
\item
  \(comp\) = Compressor
\item
  \(elec\) = Electrical
\item
  \(evap\) = Evaporator
\item
  \(grout\) = Grout material
\item
  \(high\) = High pressure side
\item
  \(in\) = Inlet condition
\item
  \(is\) = Isentropic
\item
  \(max\) = Maximum
\item
  \(motor\) = Motor
\item
  \(out\) = Outlet condition
\item
  \(pipe\) = Pipe material
\item
  \(plate\) = Plate material\\
\item
  \(pump\) = Pump
\item
  \(ref\) = Refrigerant
\item
  \(s\) = Isentropic condition
\item
  \(shaft\) = Shaft
\item
  \(soil\) = Soil/ground
\item
  \(wall\) = Wall condition
\end{itemize}

\subsubsection{Abbreviations}\label{abbreviations}

\begin{itemize}
\tightlist
\item
  \(BHE\) = Borehole Heat Exchanger
\item
  \(COP\) = Coefficient of Performance\\
\item
  \(GSHP\) = Ground Source Heat Pump
\item
  \(MEA\) = Monoethanolamine
\item
  \(OD\) = Outer Diameter
\item
  \(ID\) = Inner Diameter
\item
  \(PE\) = Polyethylene
\item
  \(SC\) = Subcooling
\item
  \(SF\) = Secondary Fluid
\item
  \(SH\) = Superheating
\end{itemize}

\subsubsection{Chemical Notation}\label{chemical-notation}

\begin{itemize}
\tightlist
\item
  \(R134A\) = 1,1,1,2-Tetrafluoroethane (refrigerant)
\item
  \(MEA\) = Monoethanolamine (\(C_2H_7NO\))
\end{itemize}

\begin{Shaded}
\begin{Highlighting}[]
\ImportTok{import} \BuiltInTok{Pkg}
\BuiltInTok{Pkg}\NormalTok{.}\FunctionTok{activate}\NormalTok{(}\StringTok{"."}\NormalTok{)}
\BuiltInTok{Pkg}\NormalTok{.}\FunctionTok{instantiate}\NormalTok{()}

\CommentTok{\# Install the local package as a regular dependency}
\BuiltInTok{Pkg}\NormalTok{.}\FunctionTok{add}\NormalTok{(path}\OperatorTok{=}\StringTok{".."}\NormalTok{)}
\end{Highlighting}
\end{Shaded}

\begin{Shaded}
\begin{Highlighting}[]
\ImportTok{using} \BuiltInTok{Revise}
\ImportTok{using} \BuiltInTok{OptimalGshpFlow}
\end{Highlighting}
\end{Shaded}

\begin{Shaded}
\begin{Highlighting}[]
\CommentTok{\# Test that the package is loaded properly}
\FunctionTok{println}\NormalTok{(}\StringTok{"OptimalGshpFlow version: "}\NormalTok{, }\FunctionTok{pkgversion}\NormalTok{(OptimalGshpFlow))}

\CommentTok{\# Create a sample property object to test functionality}
\PreprocessorTok{@named}\NormalTok{ fluid }\OperatorTok{=}\NormalTok{ Properties.}\FunctionTok{DynamicIncompressibleProps}\NormalTok{(χ }\OperatorTok{=} \FloatTok{0.2}\NormalTok{, fluid }\OperatorTok{=} \StringTok{"MEA"}\NormalTok{)}
\FunctionTok{println}\NormalTok{(}\StringTok{"Created fluid properties object: "}\NormalTok{, fluid)}
\end{Highlighting}
\end{Shaded}

\subsection{Appendix A: Model Parameters and Technical
Specifications}\label{appendix-a-model-parameters-and-technical-specifications}

\subsubsection{A.1 Borehole Heat Exchanger
Specifications}\label{a.1-borehole-heat-exchanger-specifications}

\textbf{Geometric Parameters:} - Borehole depth: \(L_b = 250\) m - Pipe
outer diameter: \(D_{pipe,o} = 32\) mm - Pipe inner diameter:
\(D_{pipe,i} = 28.24\) mm - Pipe wall thickness: \(t_{pipe} = 1.88\) mm
- Borehole diameter: \(D_{bh} = 115\) mm - Shank spacing: \(s = 75\) mm

\textbf{Thermal Properties:} - Pipe thermal conductivity:
\(\lambda_{pipe} = 0.4\) W/(m·K) - Grout thermal conductivity:
\(\lambda_{grout} = 1.5\) W/(m·K) - Ground thermal conductivity:
\(\lambda_{soil} = 2.0\) W/(m·K) - Boundary condition:
\(T_{wall} = 283.15\) K (constant)

\subsubsection{A.2 Heat Pump System
Parameters}\label{a.2-heat-pump-system-parameters}

\textbf{Working Fluids:} - Primary refrigerant: R134A - Secondary fluid:
MEA aqueous solution with mass fraction \(\chi = 0.28\)

\textbf{Compressor Specifications:} - Isentropic efficiency:
\(\eta_{is} = 0.7\) - Operating frequency: \(f = 50\) Hz - Electrical
motor efficiency: \(\eta_{motor} = 0.9\)

\subsubsection{A.3 Plate Heat Exchanger (Evaporator)
Specifications}\label{a.3-plate-heat-exchanger-evaporator-specifications}

\textbf{Geometric Parameters:} - Number of plates: \(N_{plates} = 20\) -
Plate width: \(W = 0.1\) m - Plate length: \(L = 0.3\) m - Channel
spacing: \(b = 1.5\) mm - Chevron angle: \(\beta = \pi/6\) rad \((30°)\)
- Plate thickness: \(t_{plate} = 0.5\) mm - Plate thermal conductivity:
\(\lambda_{plate} = 16\) W/(m·K)

\subsubsection{A.4 Operating Conditions (test\_GSHP Reference
Case)}\label{a.4-operating-conditions-test_gshp-reference-case}

\textbf{Temperature Differences:} - Secondary fluid temperature
difference: \(\Delta T_{SF} = 3.5\) K - Superheating temperature
difference: \(\Delta T_{SH} = 7\) K - Subcooling temperature difference:
\(\Delta T_{SC} = 5\) K

\textbf{Pressure Conditions:} - High-side pressure: \(P_{high} = 14\)
bar - Secondary fluid mass fraction: \(\chi = 0.28\)

\subsubsection{A.5 Mathematical
Correlations}\label{a.5-mathematical-correlations}

\textbf{A.5.1 Pump Performance:} - Dimensionless head coefficient:
\[\psi = \psi_0 - K \cdot \phi^2\]

\begin{itemize}
\item
  Dimensionless flow coefficient: \[\phi = \frac{\dot{V}}{N \cdot D^3}\]
\item
  Dimensionless head coefficient:
  \[\psi = \frac{g \cdot H}{N^2 \cdot D^2}\]
\item
  Pump efficiency:
  \[\eta_{pump} = \eta_{max} \left(1 - \left(\frac{\phi - \phi_{BEP}}{\phi_{BEP}}\right)^2\right)\]
\end{itemize}

where \(N\) is rotational speed {[}rpm{]}, \(D\) is impeller diameter
{[}m{]}, \(\dot{V}\) is volumetric flow rate {[}m³/s{]}, \(H\) is pump
head {[}m{]}, \(\eta_{max}\) is maximum efficiency, and \(\phi_{BEP}\)
is flow coefficient at best efficiency point.

\textbf{A.5.2 Power Consumption:} - Shaft power:
\[P_{shaft} = \frac{\rho \cdot g \cdot \dot{V} \cdot H}{\eta_{pump}}\]

\begin{itemize}
\tightlist
\item
  Electrical power: \[P_{elec} = \frac{P_{shaft}}{\eta_{motor}}\]
\end{itemize}

\textbf{A.5.3 Compressor Work:} - Isentropic compression work:
\[W_{comp} = \dot{m}_{ref} \cdot \frac{h_{out,s} - h_{in}}{\eta_{is}}\]

where \(\dot{m}_{ref}\) is refrigerant mass flow rate {[}kg/s{]},
\(h_{in}\) and \(h_{out,s}\) are inlet and isentropic outlet specific
enthalpies {[}kJ/kg{]}, respectively.

\textbf{A.5.4 Coefficient of Performance:} - Thermal COP:
\[COP_1 = \frac{\dot{Q}_{evap}}{W_{comp}}\]

\begin{itemize}
\tightlist
\item
  System COP: \[COP_2 = \frac{\dot{Q}_{evap}}{W_{comp} + P_{pump}}\]
\end{itemize}

where \(\dot{Q}_{evap}\) is evaporator heat transfer rate {[}W{]}.


\bibliography{bibliography/export.bib}



\end{document}
